\chapter{Extensión de navegador}\label{cap:extension}

\section{Motivación}

Una extensión de navegador permite al usuario invocar el sistema sin
abandonar la página que está consultando, lo que resulta especialmente
útil para verificar artículos durante la lectura.

\section{Especificación Manifest V3}

La extensión se ajusta a Manifest V3, que sustituye \emph{background
pages} por \emph{service workers} y endurece la CSP. La estructura es:

\begin{lstlisting}[caption={Layout de la extensión}]
browser-extension/
├── manifest.json
├── background/service_worker.js
├── content/{content.js, content.css}
├── popup/{popup.html, popup.css, popup.js}
├── lib/{api.js, config.js, supabase.js,
│       google-auth.js, storage.js}
└── icons/
\end{lstlisting}

\section{Autenticación}

Se reutiliza el cliente Supabase. El \emph{popup} ofrece login con
email/contraseña y con Google (vía \texttt{chrome.identity}). Las
sesiones se persisten en \texttt{chrome.storage.local} y se
sincronizan con el \emph{service worker} para inyectar el JWT en cada
llamada a la API.

\section{Capa de API}

\path{lib/api.js} define funciones \texttt{authorizedFetch},
\texttt{analyzeText} y \texttt{verifyText}. Esta última envía
\texttt{\{ texto \}} a \texttt{POST /verify} y devuelve la
\texttt{Response} parseada.

\section{Vistas del popup}

\begin{itemize}[leftmargin=*,itemsep=0.2em]
  \item \textbf{Login / Register / Forgot}: formularios con validación
        propia.
  \item \textbf{Analyze}: campo de texto libre, botón
        \emph{Analizar}, botón \emph{Verificar} (Super Pro).
  \item \textbf{Result}: muestra el veredicto y la confianza.
  \item \textbf{Verify}: muestra veredicto global, lista de
        \emph{claims} y evidencias enlazadas.
\end{itemize}

\section{Renderizado seguro}

Toda la generación de contenido dinámico utiliza \texttt{createElement}
y \texttt{textContent} en lugar de \texttt{innerHTML}, evitando
inyección de HTML desde respuestas del servidor o desde texto del
usuario.

\section{Empaquetado}

\texttt{chrome://extensions} en modo desarrollador permite cargar la
carpeta \path{browser-extension/} sin compilación, ya que el código no
usa transpilado: ES modules nativos.
